\documentclass[12pt,a4paper]{article}  % Déclare la classe du document.

%  ------------------------------------------  Modules usuels  --------------------------------------
%  --------------------------------------------------------------------------------------------------------

\usepackage{amsmath} 	
\usepackage{amssymb} 
\usepackage{listings}         
\usepackage{numprint} 
\usepackage{pifont}

%  --------------------------------------------------------------------------------------------------------
% ----------------------  packages pour fontes et caractères spéciaux  ---------------------
%  --------------------------------------------------------------------------------------------------------

\usepackage{eurosym}   
\usepackage{amsfonts}   
\usepackage{amstext}       
\usepackage{amsfonts}	
\usepackage{mathrsfs}	 

\usepackage{pstricks,pst-all,pstricks-add,pst-tree,pst-node,pst-3dplot}

%------------------------------------------------------------------------------------------------------
%  --------------------------  Francisation et gestion des caractères sous mac --------------------------------------------------------------------------------------------------------

\usepackage[francais]{babel}   
\usepackage[utf8]{inputenc}
\usepackage[OT1,T1]{fontenc} 
%  --------------------------------------------------------------------------------------------------------
% ------------------------------------Packages  de  mise en page ---------------------------
%--------------------------------------------------------------------------------------------------------

\usepackage{fancyhdr}	% en-têtes et pieds de pages
\usepackage{amsfonts}

 \usepackage{multicol} % permet l'utilisation du multicolonnage dans le texte. La syntaxe est \begin{multicols}{nb de colonnes}    ......  \end{multicols}
 
\usepackage{calc}            % permet de faire un minimum de calcul sur les arguments des commandes et macros (utilisé dans la def de l'environnement preuve )
\usepackage{epic}            % extension et simplification de  l'environnement picture
\usepackage{graphicx}    % permet d'inclure des graphique externe dans un document
\usepackage{variations}   % permet d'obtenir de joli tableau de variations
\usepackage{enumerate}  %permet de faire des listes personnalisées
\usepackage{arydshln}     %permet d'avoir des lignes en pointillés
\usepackage{framed}  
\usepackage{multido} 


%  --------------------------------------  MACROS USUELLES  ----------------------------------------

\newcommand{\R}{\mathbb{R}} 
\newcommand{\N}{\mathbb{N}} 
\newcommand{\D}{\mathbb{D}} 
\newcommand{\Z}{\mathbb{Z}} 
\newcommand{\Q}{\mathbb{Q}} 
\newcommand{\C}{\mathbb{C}} 

%  -------------------------------------  MACROS SPECIFIQUES ------------------------------------


\newcommand{\raisedrule}[2][0em]{\leaders\hbox{\rule[#1]{1pt}{#2}}\hfill}


\newcounter{numexo}

\setcounter{numexo}{1}

\newcommand{\exercice}[1]{\vskip 0.6cm\noindent\rule[+2pt]{1cm}{0.5pt}\ {\large\bf \ Exercice \thenumexo {  #1 }}\raisedrule[+2pt]{0.5pt}\null\addtocounter{numexo}{1}\par\vskip 0,4cm}


\newcommand{\lignes}[1]{
\setlength{\unitlength}{1mm}
\vskip 0,4cm\noindent
\multido{\i=0+1}{#1}{%
     \noindent
     \begin{picture}(150,5)
     \linethickness{0,005mm}
     \put(-5,0){\dashline{1}(175,5)(5,5)}
     \end{picture}\vskip 0,001cm
     }
\vskip 0cm
}




% ----------------------------------------- FORMAT DE PAGE  -----------------------------------------

\usepackage[left = 1.3cm,
	 	      right = 1.3cm,
		      top = 1.7cm,
		      headsep = 8mm,
		      bottom = 1.7cm,
		      headheight = 14.2 pt
		      ]{geometry}

\setlength{\parindent}{0mm}


\everymath{\displaystyle} 
\renewcommand{\thesection}{\Roman{section}} 
\author{}
\title{}
\date{}
\pagestyle{fancy}

 
\lhead{1ère - spécialité maths} 


\rhead{DS n$^{\circ}$6 - 26/04/21}		


\fancyfoot[C]{} 
\fancyfoot[L]{} 
\fancyfoot[R]{} 

\lstset{language=Python}  

\lstset{frame = single
framerule = 1pt
showspaces = false,
showstringspaces = false,
numbers=left, numberstyle=\tiny}

% ----------------------------------------- correction  -----------------------------------------
\usepackage[most,skins,breakable]{tcolorbox}
\usepackage{enumitem}

\newcommand{\corr}[1]{\fcolorbox{red}{yellow!30}{#1}}

\newtcolorbox{corrBox}{enhanced,boxrule=1pt,nobeforeafter,boxsep=0pt,left=\fboxsep,right=\fboxsep,colback=yellow!30,colframe=red!75!black}

\newlist{enumMet}{enumerate}{2}
\setlist[enumMet,1]{label=\arabic*), font=\bf\color{blue!75!black}}
\setlist[enumMet,2]{label=\alph*), font=\bf\color{blue!75!black}}


\begin{document}
\graphicspath{{images/}}

\begin{center}

\fbox{\LARGE{\rule{0.5cm}{0pt}\rule[-0.25cm]{0pt}{0.9cm} Devoir Surveillé N$^{\circ}$6 - Suites numériques	\rule{0.5cm}{0pt}}}
\end{center}


\begin{center}\textbf{2h - Calculatrice autorisée - Barème indicatif}\end{center}

%\begin{center}\textbf{Les élèves avec un tier-temps ne traitent pas les questions avec  le symbole \ding{46}}\end{center}

\begin{center}
%On veillera à ENCADRER ses résultats et soigner sa copie.\\
Les résultats doivent être justifiés par des calculs (au moins 1 étape intermédiaire).
\end{center}
\vspace{.05cm} 

\vspace{-0.3cm}
\exercice{ - Reconnaître une suite arithmétique - ( 3 points )}

Reconnaître parmi les suites définies sur $\mathbb{N}$ ci-dessous celles qui sont arithmétiques et préciser alors leur premier terme et leur raison.
\begin{enumMet}
    \item $u_{n}=-2+3 n$
    \item $u_{n}=2 n^{2}+3$
    \item $u_{n}=\frac{n+5}{2}$
    \item $u_{n}=1-\frac{1}{n+1}$
    \item $u_{0}=-1$ et $u_{n+1}=u_{n}+2 n$
    \item $u_{0}=4$ et $u_{n+1}=1,5+u_{n}$
\end{enumMet}

\exercice{ - Reconnaître une suite géométrique - ( 2 points )}
Reconnaître parmi les suites définies sur $\mathbb{N}$ ci-dessous celles qui sont géométriques et préciser alors leur premier terme et leur raison.
\begin{enumMet}
    \item $u_{n}=4+n \times 4$
    \item $u_{n}=\frac{2^{n}}{3}$
    \item $u_{n}=(\sqrt{2})^{n}$
    \item $u_{1}=100$ et $u_{n+1}=\frac{u_{n}}{5}$
\end{enumMet}

\exercice{ - Calculer des sommes - ( 4 points )}

\begin{enumMet}
    \item $S = 1+2+3+\ldots+1000$
    \item $T = 1+1,5+1,5^{2}+\ldots+1,5^{8}$
    \item $Z = 501+502+503+\ldots+1000$
    \item $R = 1-2+2^2+\ldots+(-2)^7$
\end{enumMet}


\exercice{ - retrouver raison et premier terme - ( 2 points )}
\begin{enumMet}
    \item On a $u_{7}=4$ et $u_{4}=7$. Calculer $r$ puis $u_{35}$
    \item On a $u_{9}=33$ et $u_{4}=13$. Calculer $r$ puis $u_{25}$.
\end{enumMet}

\exercice{ - retrouver raison et premier terme - ( 2 points )}
\begin{enumMet}
    \item On a $u_{7}=11$ et $u_{10}=3$ 773. Calculer $q$ puis $u_{12}$.
    \item On a $u_{5}=41$ et $u_{9}=25$ 625. Calculer $q$ puis $u_{10}$.
\end{enumMet}


\exercice{ - sens de variation - ( 3 points )}
Déterminer le sens de variation des suites définies pour tout entier naturel $n$ par les formules explicites suivantes.
\begin{enumMet}
    \item $u_{n}=2-4 n$
    \item $v_{n}=2 n^{2}+3$
    \item $w_{n}=n^{2}+2 n$
    \item $t_{n}=2^{n}+3^{n}$
    \item $x_{n}=n^{2}-8 n+2$
\end{enumMet}

\exercice{ - sens de variation d'une suite arithmétique - ( 2 points )}
Pour les suites arithmétiques suivantes dont on donne $\mathrm{l} \mathrm{}1^{\textrm{er}}$ terme et la raison, déterminer le sens de variation.
\begin{enumMet}
    \item $u_{0}=-2$ et $r=0,6$
    \item $v_{0}=1$ et $r=\frac{2}{3}$.
    \item $w_{0}=5$ et $r=1-\sqrt{2}$.
    \item $t_{0}=-10$ et $r=10^{-2}$.
\end{enumMet}

\exercice{ - sens de variation d'une suite arithmétique - ( 2 points )}
Pour les suites géométriques suivantes dont on donne le $1^{\text {er }}$ terme et la raison, déterminer le sens de variation
\begin{enumMet}
    \item $u_{0}=3$ et $q=2$.
    \item $v_{0}=-1$ et $q=\frac{4}{5}$.
    \item $w_{0}=\frac{-2}{3}$ et $q=\frac{8}{3}$.
    \item $t_{0}=0,5$ et $q=10^{-1}$.
\end{enumMet}

\end{document}

